\documentclass[a4paper,12pt]{article}
\usepackage[T2A]{fontenc}
\usepackage[utf8]{inputenc}
\usepackage[russian]{babel}
\usepackage{amsmath,amssymb,mathtools}
\usepackage{helvet}
\renewcommand{\familydefault}{\sfdefault}
\usepackage{icomma}
\usepackage[a4paper,left=22mm,right=22mm,top=23mm,bottom=24mm]{geometry}
\usepackage{microtype}
\usepackage{tikz}
\usetikzlibrary{arrows.meta,positioning,shapes.geometric,calc}
\usepackage{booktabs,tabularx}
\usepackage{enumitem}
\usepackage{setspace}
\usepackage{xcolor}
\usepackage{fancyhdr}
\usepackage{hyperref}
\definecolor{accent}{HTML}{008080}
\definecolor{darkaccent}{HTML}{00575A}
\definecolor{ink}{HTML}{2A3132}
\definecolor{muted}{HTML}{667174}
\hypersetup{colorlinks=true,linkcolor=darkaccent,pdfauthor={Лёвин Артём Александрович},pdftitle={Чек-лист. Тригонометрические и смешанные уравнения. 27.10.26}}
\newcommand{\key}[1]{\textcolor{darkaccent}{\textbf{#1}}}
\newcommand{\sub}[1]{\vspace{.45em}\noindent\textcolor{darkaccent}{\bfseries #1}\par\smallskip}
\setlength{\parindent}{0pt}
\setlength{\parskip}{5pt}
\setlist[enumerate]{leftmargin=8mm,itemsep=5pt,topsep=5pt}
\setstretch{1.20}
\setlength{\abovedisplayskip}{6pt}
\setlength{\belowdisplayskip}{7pt}
\setlength{\abovedisplayshortskip}{4pt}
\setlength{\belowdisplayshortskip}{5pt}
\pagestyle{fancy}
\fancyhf{}
\fancyhead[L]{\footnotesize\textcolor{muted}{Тригонометрические и смешанные уравнения}}
\fancyhead[R]{\footnotesize\textcolor{muted}{Тимофей Васильченко}}
\fancyfoot[C]{\footnotesize\textcolor{muted}{\thepage}}
\renewcommand{\headrulewidth}{0pt}
\renewcommand{\footrulewidth}{0pt}
\setlength{\headheight}{14pt}
\newcommand{\sect}[1]{\section*{\color{darkaccent}#1}}
\begin{document}
\thispagestyle{empty}
\begin{center}
\vspace*{18mm}
{\large\textcolor{darkaccent}{МАТЕМАТИКА}}\par
\vspace{18mm}
{\fontsize{23}{29}\selectfont Чек-лист}\par
\vspace{12mm}
{\fontsize{20}{25}\selectfont\textcolor{accent}{Тригонометрические}}\par
{\fontsize{20}{25}\selectfont\textcolor{accent}{и смешанные уравнения}}\par
\vspace{17mm}
{\large Подготовка к ЕГЭ по профильной математике}\par
\vfill
{\large Тимофей Васильченко}\par
\vspace{7mm}
{\large 27 октября 2026 г.}\par
\vspace{20mm}
{\small Лёвин Артём Александрович}\par
{\small эксклюзивно для Тимофея Васильченко}\par
\vspace{18mm}
\begin{tikzpicture}[x=1cm,y=1cm]
\draw[accent,line width=.75pt] (-4,0) .. controls (-2,0.85) and (-1.3,-.55) .. (0,0)
.. controls (1.6,.63) and (2.6,-.65) .. (4,0);
\fill[accent] (0,0) circle (1.5pt);
\end{tikzpicture}
\end{center}
\clearpage

\sect{1. От логарифма к тригонометрии}
\key{Главная идея.} Сначала освободитесь от логарифма с учётом его области определения. Затем применяйте тригонометрические формулы.

Для $\log_a A(x)=b$ необходимы $a>0$, $a\ne1$ и $A(x)>0$. При этих условиях
\[
\log_a A(x)=b\quad\Longleftrightarrow\quad A(x)=a^b.
\]
Уравнение решается вместе с условием $A(x)>0$. При переходе к $A(x)=a^b$ правая часть положительна, поэтому найденные корни автоматически удовлетворяют этому условию; проверка остаётся полезной.

\sub{Формулы, использованные на занятии}
\[
\begin{aligned}
\sin(\alpha-\beta)&=\sin\alpha\cos\beta-\cos\alpha\sin\beta,\\
\sin\left(x-\frac\pi2\right)&=-\cos x,\\
\sin\left(x-\frac\pi4\right)&=\frac{\sqrt2}{2}\sin x-\frac{\sqrt2}{2}\cos x,\\
\left(\frac12\right)^{1/2}&=\frac1{\sqrt2}=\frac{\sqrt2}{2}.
\end{aligned}
\]
Последняя строка использует дробный показатель степени и избавление от иррациональности в знаменателе.

\sub{Пример: логарифм с тригонометрическим аргументом}
\[
\log_{1/2}\!\left(\sin\left(x-\frac\pi4\right)
-\sqrt2\sin\left(x-\frac\pi2\right)\right)=\frac12.
\]
Пусть выражение в скобках равно $A(x)$. Записываем $A(x)>0$ и переходим к уравнению
\[
A(x)=\left(\frac12\right)^{1/2}=\frac{\sqrt2}{2}.
\]
Раскрываем оба синуса и учитываем знак перед вторым слагаемым:
\[
\frac{\sqrt2}{2}\sin x-\frac{\sqrt2}{2}\cos x
+\sqrt2\cos x=\frac{\sqrt2}{2}.
\]
Делим на ненулевое число $\sqrt2/2$ и приводим подобные:
\[
\sin x-\cos x+2\cos x=1
\quad\Longleftrightarrow\quad\boxed{\sin x+\cos x=1}.
\]
Это тригонометрическое уравнение решим методом половинного угла.

\clearpage
\sect{2. Метод половинного угла}
Если в уравнении присутствуют $\sin x$, $\cos x$ и число, удобно выразить все три части через $x/2$:
\[
\begin{aligned}
\sin x&=2\sin\frac x2\cos\frac x2,\\
\cos x&=\cos^2\frac x2-\sin^2\frac x2,\\
1&=\sin^2\frac x2+\cos^2\frac x2.
\end{aligned}
\]
\sub{Продолжение примера}
\[
\begin{aligned}
\sin x+\cos x&=1,\\
2\sin\frac x2\cos\frac x2+\cos^2\frac x2-\sin^2\frac x2
&=\sin^2\frac x2+\cos^2\frac x2,\\
2\sin\frac x2\cos\frac x2-2\sin^2\frac x2&=0,\\
2\sin\frac x2\left(\cos\frac x2-\sin\frac x2\right)&=0.
\end{aligned}
\]
Произведение равно нулю, если хотя бы один множитель равен нулю.

\textbf{Первый случай:}
\[
\sin\frac x2=0\ \Longrightarrow\ \frac x2=\pi n
\ \Longrightarrow\ \boxed{x=2\pi n}.
\]
\textbf{Второй случай:}
\[
\cos\frac x2-\sin\frac x2=0.
\]
Здесь $\cos(x/2)\ne0$: иначе $\sin(x/2)=\pm1$ и равенство было бы невозможным. Поэтому можно поделить на $\cos(x/2)$:
\[
1-\tg\frac x2=0
\ \Longrightarrow\ \tg\frac x2=1
\ \Longrightarrow\ \boxed{x=\frac\pi2+2\pi n}.
\]
В обоих семействах $n\in\mathbb Z$. При полученных значениях исходный логарифмический аргумент равен $\sqrt2/2>0$, следовательно, все корни допустимы.

\key{Ответ к примеру:}
$\displaystyle x=2\pi n$ или
$\displaystyle x=\frac\pi2+2\pi n$, $n\in\mathbb Z$.

\sub{Как распознать однородное уравнение первой степени}
После преобразований может появиться уравнение $a\sin u+b\cos u=0$. При $a\ne0$ косинус $u$ не может обращаться в нуль у решения; после проверки делим на $\cos u$ и получаем $a\tg u+b=0$.

\clearpage
\sect{3. Однородное уравнение второй степени}
Для уравнения $A\sin2x+B\cos2x=C$ раскрываем двойные углы и представляем число $C$ как $C(\sin^2x+\cos^2x)$. Тогда получим сумму слагаемых вида
\[
p\sin^2x+q\sin x\cos x+r\cos^2x=0.
\]
При $\cos x\ne0$ разделим на $\cos^2x$ и введём $t=\tg x$. Получится квадратное уравнение $pt^2+qt+r=0$.

\sub{Пример: аккуратное приведение подобных}
\[
24\sin2x-7\cos2x=15.
\]
Раскрываем обе формулы двойного угла:
\[
48\sin x\cos x-7(\cos^2x-\sin^2x)
=15(\cos^2x+\sin^2x).
\]
Переносим все слагаемые в левую часть:
\[
48\sin x\cos x-7\cos^2x+7\sin^2x
-15\cos^2x-15\sin^2x=0.
\]
Собираем одинаковые квадраты:
\[
\boxed{48\sin x\cos x-22\cos^2x-8\sin^2x=0.}
\]
\textcolor{darkaccent}{\textbf{Проверьте арифметику:}} коэффициент при $\cos^2x$ равен $-7-15=-22$.

\textbf{Нулевой косинус:} если $\cos x=0$, то $\sin2x=0$, $\cos2x=-1$, и исходное равенство становится $7=15$. Следовательно, такие значения корнями не являются.

При $\cos x\ne0$ делим на $\cos^2x$:
\[
48\tg x-22-8\tg^2x=0.
\]
После деления на $-2$ и замены $t=\tg x$:
\[
4t^2-24t+11=0,
\qquad D=24^2-4\cdot4\cdot11=400.
\]
\[
t=\frac{24\pm20}{8}
\quad\Longrightarrow\quad t=\frac12\quad\text{или}\quad t=\frac{11}{2}.
\]
\key{Ответ:}
\[
\boxed{\begin{gathered}
 x=\arctg\frac12+\pi n\quad\text{или}\\
 x=\arctg\frac{11}{2}+\pi n,\qquad n\in\mathbb Z.
\end{gathered}}
\]

\clearpage
\begin{center}
{\Large\color{darkaccent}\bfseries Алгоритм: уравнение через $\sin2x$ и $\cos2x$}
\end{center}
\vspace{14mm}
\begin{center}
\begin{tikzpicture}[
  >=Stealth,
  node distance=12mm and 15mm,
  actionbox/.style={rectangle,rounded corners=2pt,draw=accent,line width=.7pt,align=center,text width=10.1cm,minimum height=11mm,inner sep=5pt,font=\small},
  decision/.style={diamond,aspect=3,draw=accent,line width=.7pt,align=center,text width=5.8cm,inner sep=3pt,font=\small},
  side/.style={rectangle,rounded corners=2pt,draw=accent,line width=.7pt,align=center,text width=3cm,minimum height=15mm,inner sep=4pt,font=\small},
  arr/.style={->,line width=.75pt,draw=darkaccent}
]
\node[actionbox] (start) {$A\sin2x+B\cos2x=C$};
\node[actionbox,below=of start] (expand) {Замените $\sin2x$, $\cos2x$ и $1$\\формулами через $\sin x$ и $\cos x$};
\node[actionbox,below=of expand] (collect) {Соберите $\sin^2x$, $\sin x\cos x$, $\cos^2x$};
\node[decision,below=16mm of collect] (zero) {Подходит ли $\cos x=0$?};
\node[side,right=21mm of zero] (save) {Да: сохраните\\ найденные корни};
\node[actionbox,below=20mm of zero] (divide) {Для оставшихся случаев $\cos x\ne0$:\\ поделите на $\cos^2x$};
\node[actionbox,below=of divide] (quad) {Положите $t=\tg x$ и решите квадратное уравнение};
\node[actionbox,below=of quad] (result) {Восстановите $x$ и объедините все допустимые корни};
\draw[arr] (start)--(expand);
\draw[arr] (expand)--(collect);
\draw[arr] (collect)--(zero);
\draw[arr] (zero)--node[right,font=\footnotesize]{нет}(divide);
\draw[arr] (zero)--node[above,font=\footnotesize]{да}(save);
\draw[arr] (save.south) |- ([xshift=3mm]result.east);
\draw[arr] (divide)--(quad);
\draw[arr] (quad)--(result);
\end{tikzpicture}
\end{center}
\clearpage

\sect{4. Показательные уравнения с тригонометрической заменой}
\key{Главная идея.} Когда один и тот же тригонометрический фрагмент многократно входит в показатели степеней, замените его буквой и сначала решите показательное уравнение.

Если $t=\cos^2 2x$, то обязательно $\boxed{0\le t\le1}$. Числа $0$ и $1$ входят в этот диапазон.

Используйте свойства степеней:
\[
a^{m+n}=a^m a^n,
\qquad \frac{a^t}{b^t}=\left(\frac ab\right)^t
\quad (a,b>0).
\]
Если основания равны и отличны от $1$, можно приравнять показатели.

\sub{Пример: тригонометрия внутри показателей}
\[
4^{2+\cos^2 2x}-4^{\cos^2 2x}
=9^{1+\cos^2 2x}+9^{\cos^2 2x}.
\]
Обозначим $t=\cos^2 2x$, $0\le t\le1$. Тогда
\[
\begin{aligned}
4^{2+t}-4^t&=9^{1+t}+9^t,\\
16\cdot4^t-4^t&=9\cdot9^t+9^t,\\
15\cdot4^t&=10\cdot9^t,\\
\left(\frac49\right)^t&=\frac23,\\
\left(\frac23\right)^{2t}&=\left(\frac23\right)^1.
\end{aligned}
\]
Следовательно, $2t=1$, то есть $t=1/2$. Оно входит в разрешённый диапазон $[0,1]$.

Возвращаем замену и применяем формулу понижения степени:
\[
\cos^2 2x=\frac12,
\qquad \cos^2\alpha=\frac{1+\cos2\alpha}{2},
\]
\[
\frac{1+\cos4x}{2}=\frac12
\Longrightarrow \cos4x=0
\Longrightarrow 4x=\frac\pi2+\pi n.
\]
\key{Ответ:}
$\displaystyle \boxed{x=\frac\pi8+\frac{\pi n}{4},\quad n\in\mathbb Z.}$

\sub{Памятка о типичных ошибках}
\begin{enumerate}[label=\textcolor{accent}{\arabic*.},leftmargin=8mm,itemsep=3pt]
\item При раскрытии $\sin(x-\pi/2)$ сохраняйте минус перед $\cos x$.
\item Перед делением на $\cos u$ или $\cos^2u$ рассмотрите случай $\cos u=0$.
\item В формулах понижения степени для $\sin^2u$ стоит минус: $\sin^2u=(1-\cos2u)/2$.
\item После замены $t=\cos^2u$ отбрасывайте значения вне $[0,1]$.
\end{enumerate}
\clearpage

\sect{Тренировочный блок}
\textcolor{muted}{Путь А \quad\textbullet\quad 10 заданий \quad\textbullet\quad от простого к сложному}

\begin{enumerate}[label=\textcolor{darkaccent}{\bfseries\arabic*.},leftmargin=9mm,itemsep=12pt]
\item Упростите выражение $\displaystyle \sin\left(x-\frac\pi2\right)$.
\item Вычислите и избавьтесь от корня в знаменателе: $\displaystyle 2^{-1/2}$.
\item Решите уравнение $\displaystyle \cos x-\sin x=0$.
\item Решите уравнение $\displaystyle \sin x+\cos x=1$ методом половинного угла.
\item Решите уравнение $\displaystyle \sin x-\cos x=1$ методом половинного угла.
\item Решите уравнение
\[
\log_{1/2}\!\left(\sin\left(x+\frac\pi4\right)
+\sqrt2\sin\left(x-\frac\pi2\right)\right)=\frac12.
\]
\item Решите уравнение $\displaystyle 7\sin2x-\cos2x=5$.
\item Решите уравнение $\displaystyle \sin2x+3\cos2x=-1$.
\item Решите уравнение $\displaystyle \cos^2 3x=\frac12$.
\item Решите уравнение
\[
4^{1+\cos^2x}+2\cdot4^{\cos^2x}
=9^{1+\cos^2x}-5\cdot9^{\cos^2x}.
\]
\end{enumerate}

\textcolor{muted}{Во всех тригонометрических ответах $n\in\mathbb Z$. В заданиях с делением записывайте ограничения.}
\clearpage

\sect{Ответы}
\vspace{5mm}
\renewcommand{\arraystretch}{1.55}
\begin{tabularx}{\linewidth}{@{}>{\bfseries}p{12mm}X@{}}
\toprule
\textcolor{darkaccent}{№} & \textcolor{darkaccent}{Ответ} \\
\midrule
1 & $-\cos x$ \\
2 & $\displaystyle\frac{\sqrt2}{2}$ \\
3 & $\displaystyle x=\frac\pi4+\pi n$ \\
4 & $\displaystyle x=2\pi n$ или $\displaystyle x=\frac\pi2+2\pi n$ \\
5 & $\displaystyle x=\frac\pi2+2\pi n$ или $\displaystyle x=\pi+2\pi n$ \\
6 & $\displaystyle x=\frac\pi2+2\pi n$ или $\displaystyle x=\pi+2\pi n$ \\
7 & $\displaystyle x=\arctg\frac12+\pi n$ или $\displaystyle x=\arctg3+\pi n$ \\
8 & $\displaystyle x=-\frac\pi4+\pi n$ или $\displaystyle x=\arctg2+\pi n$ \\
9 & $\displaystyle x=\frac\pi{12}+\frac{\pi n}{6}$ \\
10 & $\displaystyle x=\frac\pi4+\frac{\pi n}{2}$ \\
\bottomrule
\end{tabularx}

\vspace{12mm}
\sub{Итоговый маршрут}
\begin{tabularx}{\linewidth}{@{}p{55mm}X@{}}
\toprule
\textbf{Что увидели в задаче} & \textbf{Как действовать} \\
\midrule
Логарифм от тригонометрии & Записать условие положительности, раскрыть логарифм, упростить аргумент. \\
$\sin x$ и $\cos x$ равны числу & Перейти к формулам половинного угла. \\
$\sin2x$, $\cos2x$, число & Раскрыть двойной угол, проверить $\cos x=0$, получить квадратное уравнение относительно $\tg x$. \\
Степени с $\cos^2kx$ & Ввести $t\in[0,1]$, решить показательное уравнение и вернуться к $x$. \\
\bottomrule
\end{tabularx}
\end{document}
